% TOP_PROBLEM: {"problemId": "problem.basing-one-way-functions-on-np-hardness", "canonicalTitle": "Basing One-Way Functions on NP-Hardness", "releaseRank": 33, "primaryDomain": "cryptography_coding_information_optimization", "primaryDomainLabel": "Cryptography, coding, information & optimization", "displayStatus": "open", "releaseStatus": "open"}
% !TeX program = lualatex
% Definition and sources only; this file is not an LLM solution attempt.
\documentclass[11pt]{article}
\usepackage[margin=25mm]{geometry}
\usepackage{fontspec}
\setmainfont{Latin Modern Roman}
\usepackage{amsmath,amssymb,mathtools}
\usepackage{unicode-math}
\setmathfont{Latin Modern Math}
\usepackage{xurl,hyperref}
\hypersetup{colorlinks=true,urlcolor=blue}
\setcounter{secnumdepth}{0}
\setlength{\parindent}{0pt}
\setlength{\parskip}{5pt}
\setlength{\emergencystretch}{3em}
\begin{document}
\section*{\texorpdfstring{Basing One-Way Functions on NP-Hardness}{Basing One-Way Functions on NP-Hardness}}\label{33}
\subsection{Definitions and mathematical statement}
A worst-case hardness assertion excludes an efficient algorithm correct on every input; average-case inversion hardness bounds success on $f(U_n)$, where $U_n$ is uniform on $\{0,1\}^n$. The desired construction and reduction would establish
\[
 \mathsf P\ne\mathsf{NP}\quad\Longrightarrow\quad
 \exists\text{ a one-way function family}.
\]
More operationally, an efficient inverter with nonnegligible success would be transformed into an efficient solver for an NP-hard decision problem. The class of allowed reductions matters: an obstruction to a restricted black-box reduction does not by itself refute the unrestricted goal. The catalogue leaves this reduction model broad.
\par\smallskip\textit{Formulation and concept references:} \href{https://api.openalex.org/works/doi:10.1145/1132516.1132614}{[S1]}.

\subsection{Short English statement}
Can worst-case NP-hardness alone be turned into the average-case inversion hardness needed for one-way functions?
\subsection{Sources}
\begin{itemize}
\item[S1]A. Akavia, O. Goldreich, S. Goldwasser and D. Moshkovitz, 'On basing one-way functions on NP-hardness', Proc. 38th Annual ACM Symp. on Theory of Computing (STOC), 2006, pp. 701-710. DOI:10.1145/1132516.1132614.
\url{https://api.openalex.org/works/doi:10.1145/1132516.1132614}.
\textit{Formulation source.}
\item[S2]Non-Trivial Zero-Knowledge Implies One-Way Functions.
\url{https://arxiv.org/abs/2602.17651}.
\textit{Further reference; not a proof endorsement.}
\item[S3]Hardness Along the Boundary: Towards One-Way Functions from the Worst-case Hardness of Time-Bounded Kolmogorov Complexity.
\url{https://eccc.weizmann.ac.il/report/2025/125/}.
\textit{Further reference; not a proof endorsement.}
\item[S4]Capturing One-Way Functions via NP-Hardness of Meta-Complexity.
\url{https://eccc.weizmann.ac.il/report/2023/037/}.
\textit{Further reference; not a proof endorsement.}
\end{itemize}
\end{document}
